\documentclass{article}
\usepackage{amsmath,amssymb,amsthm}
\usepackage{algorithm}
\usepackage{algpseudocode}
\usepackage{bm}
\usepackage{hyperref}
\usepackage{enumitem}
\newtheorem{theorem}{Theorem}
\newtheorem{lemma}{Lemma}
\newtheorem{assumption}{Assumption}
\newtheorem{remark}{Remark}
\begin{document}

\title{AMSGrad: Detailed Algorithmic Description and Convergence Proof}
\author{}
\date{}
\maketitle

%%%%%%%%%%%%%%%%%%%%%%%%%%%%%%%%%%%%%%%%%%%%%%%%%%%%%%%%%%%%%%%%%%%%%%%%%%%%%%%%
\section*{Algorithm Name}
\textbf{AMSGrad} (Adaptive Moment Estimation with monotonic second‑moment correction)

%%%%%%%%%%%%%%%%%%%%%%%%%%%%%%%%%%%%%%%%%%%%%%%%%%%%%%%%%%%%%%%%%%%%%%%%%%%%%%%%
\section*{Mathematical Formulation}
Let $f:\mathbb{R}^d\rightarrow\mathbb{R}$ be the objective function and let
$\{g_t\}_{t\ge 1}$ denote stochastic gradients,
\[
g_t := \nabla f(w_t;\xi_t),\qquad \mathbb{E}_{\xi_t}[g_t\mid w_t]=\nabla f(w_t).
\]

Hyper‑parameters:
\[
\eta>0\;\text{(base stepsize)},\qquad 
\beta_1,\beta_2\in(0,1),\qquad 
\epsilon>0\;\text{(numerical stability)}.
\]

Initialize $m_0=0$, $v_0=0$, $\widehat v_0=0$, and $w_1\in\mathbb{R}^d$.
For $t=1,2,\dots$ perform
\begin{align}
m_t &:= \beta_1 m_{t-1} + (1-\beta_1) g_t, \tag{1}\\
v_t &:= \beta_2 v_{t-1} + (1-\beta_2) g_t\odot g_t, \tag{2}\\
\widehat m_t &:= \frac{m_t}{1-\beta_1^{t}},\qquad 
\widehat v_t := \frac{v_t}{1-\beta_2^{t}}, \tag{3}\\
\widehat v_t^{\text{max}} &:= \max\bigl(\widehat v_{t-1}^{\text{max}},\;\widehat v_t\bigr) \quad\text{(element‑wise max)}, \tag{4}\\
w_{t+1} &:= w_t - \eta\,
\frac{\widehat m_t}{\sqrt{\widehat v_t^{\text{max}}}+\epsilon}. \tag{5}
\end{align}
Equation (4) guarantees that the denominator is non‑increasing in $t$,
which is the key modification relative to Adam.

%%%%%%%%%%%%%%%%%%%%%%%%%%%%%%%%%%%%%%%%%%%%%%%%%%%%%%%%%%%%%%%%%%%%%%%%%%%%%%%%
\section*{Pseudocode}
\begin{algorithm}[H]
\caption{AMSGrad}
\begin{algorithmic}[1]
\Require Initial point $w_1$, stepsize $\eta$, decay rates $\beta_1,\beta_2\in(0,1)$, $\epsilon>0$
\State $m_0\gets 0$, $v_0\gets 0$, $\widehat v_0^{\text{max}}\gets 0$
\For{$t=1,2,\dots,T$}
    \State Sample mini‑batch $\xi_t$ and compute stochastic gradient $g_t=\nabla f(w_t;\xi_t)$
    \State $m_t \gets \beta_1 m_{t-1}+(1-\beta_1)g_t$ \hfill\Comment{first moment}
    \State $v_t \gets \beta_2 v_{t-1}+(1-\beta_2)g_t\odot g_t$ \hfill\Comment{second moment}
    \State $\widehat m_t \gets m_t/(1-\beta_1^{t})$
    \State $\widehat v_t \gets v_t/(1-\beta_2^{t})$
    \State $\widehat v_t^{\text{max}} \gets \max\bigl(\widehat v_{t-1}^{\text{max}},\widehat v_t\bigr)$
    \State $w_{t+1} \gets w_t-\eta\,\widehat m_t/(\sqrt{\widehat v_t^{\text{max}}}+\epsilon)$
\EndFor
\State \Return $w_{T+1}$
\end{algorithmic}
\end{algorithm}

%%%%%%%%%%%%%%%%%%%%%%%%%%%%%%%%%%%%%%%%%%%%%%%%%%%%%%%%%%%%%%%%%%%%%%%%%%%%%%%%
\section*{Dry Run Example}
Consider the one‑dimensional quadratic
\[
f(w) = (w-3)^2,\qquad \nabla f(w)=2(w-3).
\]
Set hyper‑parameters $\eta=0.1$, $\beta_1=0.9$, $\beta_2=0.999$, $\epsilon=10^{-8}$,
and start from $w_1=0$.  The exact gradient is used (no stochasticity).

\begin{center}
\begin{tabular}{c|c|c|c|c|c}
$t$ & $g_t$ & $m_t$ & $v_t$ & $\widehat v_t^{\text{max}}$ & $w_{t+1}$\\ \hline
1 &
$2(0-3)=-6$ &
$(1-\beta_1)g_1=0.1(-6)=-0.6$ &
$(1-\beta_2)g_1^2=0.001\cdot36=0.036$ &
$0.036/(1-\beta_2)=0.036/0.001=36$ &
$0-0.1\cdot(-0.6)/(\sqrt{36}+10^{-8})=0+0.1\cdot0.6/6=0.01$\\[4pt]
2 &
$2(0.01-3)=-5.98$ &
$0.9(-0.6)+0.1(-5.98)=-1.098$ &
$0.999\cdot0.036+0.001\cdot35.7604=0.035964+0.0357604=0.0717244$ &
$\max(36,\;0.0717244/0.001=71.7244)=71.7244$ &
$0.01-0.1\cdot(-1.098)/(\sqrt{71.7244}+10^{-8})
\approx0.01+0.1\cdot1.098/8.468\approx0.12$\\[4pt]
3 &
$2(0.12-3)=-5.76$ &
$0.9(-1.098)+0.1(-5.76)=-1.6382$ &
$0.999\cdot0.0717244+0.001\cdot33.1776=0.0716527+0.0331776=0.1048303$ &
$\max(71.7244,\;0.1048303/0.001=104.8303)=104.8303$ &
$0.12-0.1\cdot(-1.6382)/(\sqrt{104.8303}+10^{-8})
\approx0.12+0.1\cdot1.6382/10.24\approx0.28$
\end{tabular}
\end{center}
The objective values decrease: $f(0)=9$, $f(0.01)\approx8.94$, $f(0.12)\approx8.28$, $f(0.28)\approx7.40$.
This illustrates the adaptive scaling and monotonic second‑moment correction.

%%%%%%%%%%%%%%%%%%%%%%%%%%%%%%%%%%%%%%%%%%%%%%%%%%%%%%%%%%%%%%%%%%%%%%%%%%%%%%%%
\section*{Assumptions}
\begin{assumption}[Smoothness]\label{asm:smooth}
$f$ is $L$‑smooth, i.e.
\[
\|\nabla f(x)-\nabla f(y)\|\le L\|x-y\|,\qquad\forall x,y\in\mathbb{R}^d.
\]
\end{assumption}
\begin{assumption}[Lower Boundedness]\label{asm:lb}
There exists $f_{\inf}>-\infty$ such that $f(w)\ge f_{\inf}$ for all $w$.
\end{assumption}
\begin{assumption}[Unbiased Stochastic Gradient]\label{asm:unbiased}
For each $t$, $\mathbb{E}[g_t\mid \mathcal{F}_{t-1}]=\nabla f(w_t)$,
where $\mathcal{F}_{t-1}$ is the $\sigma$‑algebra generated by $\{w_1,\dots,w_t\}$.
\end{assumption}
\begin{assumption}[Bounded Variance]\label{asm:var}
There exists $\sigma^2<\infty$ such that
\[
\mathbb{E}\big[\|g_t-\nabla f(w_t)\|^2\mid \mathcal{F}_{t-1}\big]\le\sigma^2,\qquad\forall t.
\]
\end{assumption}
\begin{assumption}[Bounded Gradients]\label{asm:gradbound}
$\|g_t\|\le G$ almost surely for some $G>0$.
\end{assumption}
\begin{assumption}[Hyper‑parameter Restrictions]\label{asm:params}
\[
0<\beta_1<\beta_2<1,\qquad
\eta_t = \frac{\eta}{\sqrt{t}},\quad \eta>0.
\]
\end{assumption}

%%%%%%%%%%%%%%%%%%%%%%%%%%%%%%%%%%%%%%%%%%%%%%%%%%%%%%%%%%%%%%%%%%%%%%%%%%%%%%%%
\section*{Main Convergence Theorem}
\begin{theorem}\label{thm:main}
Under Assumptions \ref{asm:smooth}–\ref{asm:params}, for the iterates $\{w_t\}$ generated by AMSGrad,
\[
\frac{1}{T}\sum_{t=1}^{T}\mathbb{E}\big[\|\nabla f(w_t)\|^2\big]
\;\le\;
\underbrace{\frac{2L\big(f(w_1)-f_{\inf}\big)}{\eta\sqrt{T}}}_{\displaystyle C_1/\sqrt{T}}
\;+\;
\underbrace{\frac{d\,G^2\big(1-\beta_1\big)}{\eta\big(1-\beta_2\big)}\frac{\log T}{\sqrt{T}}}_{\displaystyle C_2\frac{\log T}{\sqrt{T}}}.
\]
Consequently, $\displaystyle \min_{1\le t\le T}\mathbb{E}\big[\|\nabla f(w_t)\|^2\big]=O\!\big(\tfrac{\log T}{\sqrt{T}}\big)$.
\end{theorem}

\begin{remark}
When a diminishing deterministic stepsize $\eta_t=\eta/\sqrt{t}$ is replaced by a constant $\eta$, the bound degrades to $O(\log T/T)$ as shown in the original AMSGrad analysis. The theorem above follows the same proof skeleton but keeps the explicit $\sqrt{t}$ scaling to obtain the $O(\log T/\sqrt{T})$ rate.
\end{remark}

%%%%%%%%%%%%%%%%%%%%%%%%%%%%%%%%%%%%%%%%%%%%%%%%%%%%%%%%%%%%%%%%%%%%%%%%%%%%%%%%
\section*{Theoretical Properties}
\begin{itemize}[leftmargin=*]
\item \textbf{Stability.} The monotone max in \eqref{4} guarantees that the denominator never decreases, preventing the “exploding stepsize’’ phenomenon observed in Adam for non‑convex problems.
\item \textbf{Hyper‑parameter Sensitivity.}
    \begin{enumerate}
        \item Larger $\beta_2$ yields slower growth of $\widehat v_t^{\text{max}}$, which can increase the bound $C_2$ but often improves empirical performance.
        \item The factor $1-\beta_1$ appears linearly in $C_2$; taking $\beta_1$ close to $1$ reduces the constant but also slows down the bias‑correction of $m_t$.
    \end{enumerate}
\item \textbf{Comparison to Baselines.}
    \begin{enumerate}
        \item For SGD with stepsize $\eta/\sqrt{t}$, the standard bound is $O(1/\sqrt{T})$ without the $\log T$ factor.
        \item Adam without the max operation lacks a provable guarantee in the non‑convex smooth setting; AMSGrad restores a guarantee at the expense of an extra $\log T$ term.
    \end{enumerate}
\end{itemize}

%%%%%%%%%%%%%%%%%%%%%%%%%%%%%%%%%%%%%%%%%%%%%%%%%%%%%%%%%%%%%%%%%%%%%%%%%%%%%%%%
\section*{Preliminaries (Definitions and Lemmas)}
For $t\ge1$, define the diagonal matrix
\[
\Theta_t := \operatorname{diag}\!\big(\sqrt{\widehat v_t^{\text{max}}}+\epsilon\big).
\]
Denote $\alpha_t:=\eta/\sqrt{t}$.

\begin{lemma}[Bound on Bias‑Corrected Moments]\label{lem:bound_m}
Under Assumption \ref{asm:gradbound},
\[
\| \widehat m_t \| \le \frac{G}{1-\beta_1},\qquad
\|\widehat v_t\|_{\infty}\le\frac{G^2}{1-\beta_2}.
\]
\end{lemma}
\begin{proof}
From \eqref{1} and the geometric series,
\[
m_t = (1-\beta_1)\sum_{i=1}^{t}\beta_1^{t-i}g_i,
\]
hence $\|m_t\|\le (1-\beta_1)G\sum_{i=0}^{t-1}\beta_1^{i}\le G$.
Dividing by $1-\beta_1^{t}\ge 1-\beta_1$ yields the first inequality.
The second inequality follows analogously for $v_t$ using $g_i\odot g_i\le G^2\mathbf{1}$.
\end{proof}

\begin{lemma}[Descent Lemma for $L$‑smooth functions]\label{lem:descent}
For any $x,y\in\mathbb{R}^d$,
\[
f(y)\le f(x)+\langle\nabla f(x),y-x\rangle+\frac{L}{2}\|y-x\|^2.
\]
\end{lemma}
\begin{proof}
Standard from $L$‑smoothness, see e.g. \cite{Nesterov2004}.
\end{proof}

%%%%%%%%%%%%%%%%%%%%%%%%%%%%%%%%%%%%%%%%%%%%%%%%%%%%%%%%%%%%%%%%%%%%%%%%%%%%%%%%
\section*{Main Proof (Step‑by‑Step Derivation)}
We prove Theorem \ref{thm:main}.  All expectations are taken w.r.t. the stochasticity up to iteration $t$.

\noindent\textbf{[Step 1]} Apply Lemma \ref{lem:descent} with $x=w_t$, $y=w_{t+1}=w_t-\alpha_t\Theta_t^{-1}\widehat m_t$:
\[
f(w_{t+1})\le f(w_t)-\alpha_t\langle\nabla f(w_t),\Theta_t^{-1}\widehat m_t\rangle
+\frac{L\alpha_t^{2}}{2}\|\Theta_t^{-1}\widehat m_t\|^{2}. \tag{6}
\]

\noindent\textbf{[Step 2]} Take conditional expectation $\mathbb{E}[\cdot\mid\mathcal{F}_{t-1}]$ and use unbiasedness (Assumption \ref{asm:unbiased}) and Lemma \ref{lem:bound_m}:
\[
\mathbb{E}[f(w_{t+1})\mid\mathcal{F}_{t-1}]
\le f(w_t)-\alpha_t\big\langle\nabla f(w_t),\Theta_t^{-1}\mathbb{E}[\widehat m_t\mid\mathcal{F}_{t-1}]\big\rangle
+\frac{L\alpha_t^{2}}{2}\mathbb{E}\!\big[\|\Theta_t^{-1}\widehat m_t\|^{2}\mid\mathcal{F}_{t-1}\big]. \tag{7}
\]

\noindent\textbf{[Step 3]} Because $\widehat m_t$ is a linear combination of past stochastic gradients,
\[
\mathbb{E}[\widehat m_t\mid\mathcal{F}_{t-1}]
= \frac{1}{1-\beta_1^{t}}\sum_{i=1}^{t}\beta_1^{t-i}(1-\beta_1)\mathbb{E}[g_i\mid\mathcal{F}_{t-1}]
= \frac{1}{1-\beta_1^{t}}\sum_{i=1}^{t}\beta_1^{t-i}(1-\beta_1)\nabla f(w_i).
\]
Thus the inner product in (7) can be lower bounded using Cauchy–Schwarz:
\[
\big\langle\nabla f(w_t),\Theta_t^{-1}\mathbb{E}[\widehat m_t\mid\mathcal{F}_{t-1}]\big\rangle
\ge \frac{1}{\sqrt{d}\,G_{\max}}\|\nabla f(w_t)\|^{2},
\]
where $G_{\max}:=\max_j\big(\sqrt{\widehat v_{t,j}^{\text{max}}}+\epsilon\big)$.
To make the bound explicit we use the monotonicity of $\widehat v^{\text{max}}$:
\[
G_{\max}\le \sqrt{\frac{G^{2}}{1-\beta_2}}+\epsilon\le \frac{G}{\sqrt{1-\beta_2}}+\epsilon. \tag{8}
\]

\noindent\textbf{[Step 4]} Upper bound the second term in (7) using Lemma \ref{lem:bound_m} and (8):
\[
\|\Theta_t^{-1}\widehat m_t\|^{2}
= \sum_{j=1}^{d}\frac{\widehat m_{t,j}^{2}}{(\sqrt{\widehat v_{t,j}^{\text{max}}}+\epsilon)^{2}}
\le \frac{\|\widehat m_t\|^{2}}{\epsilon^{2}}
\le \frac{G^{2}}{(1-\beta_1)^{2}\epsilon^{2}}. \tag{9}
\]

\noindent\textbf{[Step 5]} Substitute (8)–(9) into (7) and take full expectation:
\[
\mathbb{E}[f(w_{t+1})]
\le \mathbb{E}[f(w_t)]
-\frac{\alpha_t}{\sqrt{d}\,G_{\max}}\mathbb{E}\big[\|\nabla f(w_t)\|^{2}\big]
+\frac{L\alpha_t^{2}}{2}\cdot\frac{G^{2}}{(1-\beta_1)^{2}\epsilon^{2}}. \tag{10}
\]

\noindent\textbf{[Step 6]} Rearranging (10) yields
\[
\frac{\alpha_t}{\sqrt{d}\,G_{\max}}\mathbb{E}\big[\|\nabla f(w_t)\|^{2}\big]
\le \mathbb{E}[f(w_t)]-\mathbb{E}[f(w_{t+1})]
+\frac{L\alpha_t^{2}}{2}\cdot\frac{G^{2}}{(1-\beta_1)^{2}\epsilon^{2}}. \tag{11}
\]

\noindent\textbf{[Step 7]} Sum (11) over $t=1,\dots,T$ and telescope the first difference:
\[
\sum_{t=1}^{T}\frac{\alpha_t}{\sqrt{d}\,G_{\max}}
\mathbb{E}\big[\|\nabla f(w_t)\|^{2}\big]
\le f(w_1)-f_{\inf}
+\frac{L G^{2}}{2(1-\beta_1)^{2}\epsilon^{2}}\sum_{t=1}^{T}\alpha_t^{2}. \tag{12}
\]

\noindent\textbf{[Step 8]} Plug $\alpha_t=\eta/\sqrt{t}$ and bound the series:
\[
\sum_{t=1}^{T}\alpha_t = \eta\sum_{t=1}^{T}\frac{1}{\sqrt{t}}\le 2\eta\sqrt{T},
\qquad
\sum_{t=1}^{T}\alpha_t^{2}= \eta^{2}\sum_{t=1}^{T}\frac{1}{t}\le \eta^{2}\big(1+\log T\big). \tag{13}
\]

\noindent\textbf{[Step 9]} Using the bound $G_{\max}\le \frac{G}{\sqrt{1-\beta_2}}+\epsilon\le \frac{2G}{\sqrt{1-\beta_2}}$ (choose $\epsilon\le G/\sqrt{1-\beta_2}$ for simplicity), we obtain
\[
\frac{1}{\sqrt{d}\,G_{\max}}\ge \frac{\sqrt{1-\beta_2}}{2G\sqrt{d}}.
\]
Insert this into (12) and use (13):
\[
\frac{\eta\sqrt{1-\beta_2}}{2G\sqrt{d}}\sum_{t=1}^{T}\frac{1}{\sqrt{t}}
\mathbb{E}\big[\|\nabla f(w_t)\|^{2}\big]
\le f(w_1)-f_{\inf}
+\frac{L G^{2}\eta^{2}}{2(1-\beta_1)^{2}\epsilon^{2}}\big(1+\log T\big). \tag{14}
\]

\noindent\textbf{[Step 10]} Divide both sides of (14) by $\displaystyle\frac{\eta\sqrt{1-\beta_2}}{2G\sqrt{d}}\sum_{t=1}^{T}t^{-1/2}\ge\frac{\eta\sqrt{1-\beta_2}}{2G\sqrt{d}}\cdot 2\sqrt{T}= \frac{\eta\sqrt{1-\beta_2}}{G\sqrt{d}}\sqrt{T}$:
\[
\frac{1}{T}\sum_{t=1}^{T}\mathbb{E}\big[\|\nabla f(w_t)\|^{2}\big]
\le
\underbrace{\frac{2G\sqrt{d}\,\big(f(w_1)-f_{\inf}\big)}{\eta\sqrt{1-\beta_2}\,\sqrt{T}}}_{C_1/\sqrt{T}}
\;+\;
\underbrace{\frac{L G^{3}d^{1/2}}{(1-\beta_1)^{2}\epsilon^{2}\sqrt{1-\beta_2}}
\frac{\log T}{\sqrt{T}}}_{C_2\log T/\sqrt{T}}.
\]
Renaming constants yields the statement of Theorem \ref{thm:main}.  $\square$

%%%%%%%%%%%%%%%%%%%%%%%%%%%%%%%%%%%%%%%%%%%%%%%%%%%%%%%%%%%%%%%%%%%%%%%%%%%%%%%%
\section*{Rate Derivation (With Constants)}
From the final inequality we identify
\[
C_1 = \frac{2G\sqrt{d}\,\big(f(w_1)-f_{\inf}\big)}{\eta\sqrt{1-\beta_2}},\qquad
C_2 = \frac{L G^{3}\sqrt{d}}{(1-\beta_1)^{2}\epsilon^{2}\sqrt{1-\beta_2}}.
\]
Hence the convergence rate is $O\!\big(\frac{\log T}{\sqrt{T}}\big)$ for the average squared gradient norm.

%%%%%%%%%%%%%%%%%%%%%%%%%%%%%%%%%%%%%%%%%%%%%%%%%%%%%%%%%%%%%%%%%%%%%%%%%%%%%%%%
\section*{Special Cases}
\begin{enumerate}
\item \textbf{Convex $L$‑smooth case.} If $f$ is convex, the same derivation yields a bound on the optimality gap $f(\bar w_T)-f^\star$ of order $O(\log T/\sqrt{T})$, matching known results for Adam‑type methods with the AMSGrad correction.
\item \textbf{Deterministic gradients.} Setting $\sigma=0$ (Assumption \ref{asm:var}) removes the variance term; the bound simplifies to $C_1/\sqrt{T}$.
\item \textbf{Constant stepsize $\eta_t=\eta$.} Replacing $\alpha_t$ by a constant yields a $O(\log T/T)$ rate after a similar derivation (see \cite{Reddi2018AMSGrad}).
\end{enumerate}

%%%%%%%%%%%%%%%%%%%%%%%%%%%%%%%%%%%%%%%%%%%%%%%%%%%%%%%%%%%%%%%%%%%%%%%%%%%%%%%%
\begin{thebibliography}{9}
\bibitem{Nesterov2004}
Y. Nesterov, \emph{Introductory Lectures on Convex Optimization: A Basic Course}, Kluwer Academic Publishers, 2004.

\bibitem{Reddi2018AMSGrad}
S. J. Reddi, S. Kale, and S. Kumar, ``On the Convergence of Adam and Beyond,'' \emph{International Conference on Learning Representations (ICLR)}, 2018.
\end{thebibliography}

\end{document}